\section{研究问题与共同方法：结构、状态和证据必须分开}
\label{sec:questions}
\subsection{为什么从习惯化开始}
成年果蝇连接组将大量神经元及其突触接触关系转化为可计算的结构约束\citep{dorkenwald2024}。Shiu等人的全脑泄漏积分发放模型进一步表明，这类约束能够支持感觉到脑内输出的计算预测\citep{shiu2024}。然而，正常传播与历史依赖是两种不同能力：一次输入能够激活正确读出，并不意味着反复输入会使其改变。

习惯化提供了一个相对简洁的切入口。它不要求语言任务、显式奖励或大规模权重优化，却要求当前响应受到过去刺激的影响。重复刺激后响应下降、休息后恢复，是容易测量的起点，但输入变弱、外围适应和输出失能也能产生类似表象\citep{rankin2009}。因此，本研究不以“画出一条下降曲线”为终点，而逐步追问该现象是否由已知状态造成、能否预测未见条件、以及还缺少哪些特征。

最小动力学研究已经说明，少量内部状态配合非线性输出就能产生多项习惯化特征\citep{smart2024}。这既为最小机制提供思路，也提高了连接组模型的证据要求：全脑规模本身不是解释力的证明。嗅觉和味觉研究提示抑制性传递可塑性与解除习惯化回路是合理竞争机制\citep{das2011,paranjpe2012,trisal2022}；视觉行为和个体动力学研究则提醒我们注意协议与个体差异\citep{engel2009,boon2026}。这些资料不能被直接当作JO-CE机制的证据。

\subsection{五个递进问题，而不是五份互不相干的结果}
\begin{table}[!htbp]\centering\small
\caption{统一研究问题及其证据职责。后一个问题回应前一层证据尚不能回答的内容。}
\begin{tabularx}{\linewidth}{p{.8cm}YY}\toprule
问题 & 要回答什么 & 对应检验与边界\\\midrule
Q1 & 真实连接和原始快动力学是否足以产生递减？ & M0与局部STD的M1比较；阳性只是条件性充分性。\\
Q2 & 历史储存在哪里，输出通路是否失能？ & 资源/快状态干预及旁路；只在特定模型内部归因。\\
Q3 & 小模型能否解释并预测这些输出？ & 旧数据拟合、冻结预测、新条件评分；不等于识别动物机制。\\
Q4 & 能否推广到更多习惯化特征？ & 多轮、强度、恢复、追加训练和B刺激；按资格与协议逐项解释。\\
Q5 & 真实机制和连接结构的独特价值是什么？ & 当前仍缺活体数据、局部/随机图对照及细胞干预预测。\\\bottomrule
\end{tabularx}\end{table}

这条路线改变的是证据要求，不是不断给模型增加参数。Q1后没有继续调整全脑参数使曲线更漂亮；Q2后转而检验更简单的描述；Q3后保留其不能解释的反弹；Q4则检验此前尚不能声称具备的特征。Q5仍然开放。

\subsection{模型范围、输入与读出}
全部研究使用固定FlyWire v630和同一上游提交，包含127,400个神经元、14,687,178条聚合有向边，对应52,793,639个解剖突触。多个解剖接触被聚合成一条边，三种数量不能混用。没有裁剪子图；“全脑”指该模型的数据范围，不包含虚拟身体、完整腹神经索或所有细胞生理过程。

输入是从固定Notebook清单取得的70个JO-CE机械感觉神经元，主要读出为触角梳理相关脑内中间神经元aBN1，次要读出为相关下行神经元aDN1/aDN2。输入是人工神经激活，不是已校准的声音、机械力或触角位移。aBN1计数也不是梳理概率或运动幅度。已有糖/苦味复现只作为工程背景，不混入这里的习惯化样本数量。

对每个种子，在时间步$\Delta t=0.1$ms上生成Bernoulli事件；脉冲内命令输入率为$f_{\rm cmd}$时，每格概率为$f_{\rm cmd}\Delta t$。同一种子的单次事件阵列在相应重复刺激和配对条件间回放。输入细胞不应期按接口设置为零，注入幅度沿用上游$250w_{\rm syn}=68.75$mV。命令事件经过阈值/重置调度才变成实际感觉脉冲，二者不是严格一对一；因此必须核对\emph{实际}感觉事件，不能只核对随机种子或命令计数。

\subsection{M0与M1：只引入一种候选慢机制}
对不处于不应期的神经元$j$，M0沿用
\begin{align}
\tau_m\dot v_j&=v_0-v_j+g_j, & \tau_g\dot g_j&=-g_j,\label{eq:lif}\\
g_j&\leftarrow g_j+n_{ij}s_iw_{\rm syn} &&\text{（突触前事件延迟到达时）。}
\end{align}
$g$是电压量纲的有效驱动，不是电导。越过阈值后发放，将$v$置为重置电位、$g$清零，并进入不应期；上游实现中两个微分方程在不应期均暂停更新。数值实现使用Brian2线性更新与Cython后端\citep{stimberg2019}。

\begin{table}[!htbp]\centering\small
\caption{共同参数。M1的两个共享参数未用本研究数据做生理拟合。}
\begin{tabularx}{\linewidth}{lrlY}\toprule
参数 & 数值 & 单位 & 作用\\\midrule
$v_0,v_{\rm rst}$ & $-52$ & mV & 静息/重置电位\\
$v_{\rm th}$ & $-45$ & mV & 发放阈值\\
$\tau_m,\tau_g$ & $20,5$ & ms & 膜与驱动时间常数\\
$t_{\rm rfc},t_{\rm dly}$ & $2.2,1.8$ & ms & 普通细胞不应期、传递延迟\\
$w_{\rm syn}$ & $0.275$ & mV & 解剖突触计数的全局尺度\\
$U$ & $0.01$ & 无量纲 & 单次事件后的资源消耗比例\\
$\tau_{\rm rec}$ & $2$ & s & 新增资源恢复尺度\\\bottomrule
\end{tabularx}\end{table}

M1只把JO-CE的4,764条正权重输出边替换为同拓扑、同延迟的动态边；原对应边置零，避免双重传递。资源初值$x_{ij}=1$，事件间及事件到达时分别满足
\begin{align}
\dot x_{ij}&=(1-x_{ij})/\tau_{\rm rec},\label{eq:resource}\\
g_j&\leftarrow g_j+w_{ij}x_{ij}^{-},\qquad
x_{ij}^{+}=(1-U)x_{ij}^{-}.\label{eq:event}
\end{align}
先传递、后消耗，事件间按指数精确更新。STD指\emph{兴奋性传递的短时效能下降}，不是把连接变成抑制性连接，也不是加强一群抑制性神经元。

这个$x$是可恢复资源的抽象状态，不是实测囊泡数。所有边共享$U$和$\tau_{\rm rec}$，并不意味着全脑只有一个标量状态，更不意味着4,764个参数被单独训练。在同一突触前细胞、延迟和初态下，部分边的资源历史可以相同；本研究未进行状态共享优化。

\paragraph{三个“训练”概念。}
实验中的重复刺激是暴露历史；M1的资源变化是短时动态可塑性；后文代理的最小二乘才是参数拟合。全脑解剖结构和生理参数始终不通过优化器训练。固定权重网络仍有膜电位和循环活动等状态，不能先验断言它绝不可能记忆。

\subsection{统一计量：时间、分母和统计单位}
令刺激宽度为$P$，响应窗口为$W$，起始间隔为$T$。训练序列内不重建网络；通常$W=P+100$ms。$T$不是静默时间，刺激结束到下一次起始的间隔是$T-P$。休息探测从最后\emph{响应窗口结束}计时，距刺激结束还多100ms。

训练末态可以保存后分叉，每个探测从同一个末态独立等待；这不是将资源恢复到初态。多轮训练之间则沿同一轨迹连续等待，不通过初态恢复制造“自发恢复”。人工状态救援是另一类明确标记的干预。

对一条含$N$次刺激的序列，定义
\begin{equation}
I=R_1,\quad E=\frac{R_1+R_2+R_3}{3},\quad L=\frac{R_{N-2}+R_{N-1}+R_N}{3},\quad D=1-\frac{L}{E}.
\label{eq:D}
\end{equation}
$D$先逐种子计算再平均，$E=0$时未定义。初次至少5脉冲为正常传递门槛，$D\geq0.30$为递减筛查；相邻两组三次末期均值之差不超过$0.10E$为平台筛查。这些是本项目效应量约定，不是经典文献的普适阈值；通过D筛查也不等于每一步都严格单调不增。

两个恢复读出分别为
\begin{equation}
Q(t)=\frac{R_{\rm probe}(t)}{I},\qquad F(t)=\frac{R_{\rm probe}(t)-L}{I-L},\quad I>L.
\label{eq:normalization}
\end{equation}
$Q$衡量离初次水平多远，$F$衡量已损失部分恢复多少。负值或过冲不截断，两者不能互换成恢复时间常数。

三个阶段使用1701--1705、2701--2703、3701--3703，共11个不同种子标识；因果阶段复用第一组。它们都是\textbf{一个图谱上的输入实现，不是11只动物}。序列内响应相互依赖，代理拟合窗口还与旧实验重叠。本文报告描述性计数、范围和筛查数，不据此作动物总体显著性推断。
\FloatBarrier
